\chapter*{ABSTRACT}

\begin{center}
    \textbf{\fontsize{13}{16}\selectfont CAUSAL INFERENCE AND DEEP GRAPHICAL REPRESENTATIONS FOR EDUCATIONAL POLICY INTERVENTION}
\end{center}
\vspace{0.3cm}

Estimating the true causal impact of educational interventions from observational tabular records is a fundamental problem in modern institutional policy design. While traditional Causal Bayesian Networks (CBNs) provide rigorous identification guarantees and transparent probabilistic reasoning through Judea Pearl's do-calculus, discrete conditional probability tables (CPDs) suffer from statistical sparsity and struggle to capture continuous or high-order nonlinear interactions. Conversely, purely correlational deep learning and machine learning algorithms achieve high predictive accuracy but fail to identify counterfactual outcomes due to confounding bias and collider stratification. 

In this study, we propose and rigorously benchmark a controlled hybrid framework—\textbf{Hybrid Neural-Causal Probabilistic Graphical Models (Hybrid PGM)}—that bridges structural causal discovery, Bayesian graphical modeling, and deep neural parameterization. Focusing on the public educational success benchmark (UCI Machine Learning Repository ID 697, $N=4,424$, binary cohort $N=3,630$) and its large-scale counterpart (Kaggle Playground Series S4E6, $N=76,518$), we formalize a 7-node Primary Domain Directed Acyclic Graph (DAG) with 11 directed edges, separating 24 pre-enrollment baseline covariates from 12 post-enrollment academic progression variables. We design a $2 \times 2$ factorial evaluation comparing four paradigms: baseline tabular CBN with BDeu smoothing ($B_0$), fixed-DAG hybrid neural outcome parameterization via PyTorch Multi-Layer Perceptrons ($H_1$), continuous optimization-based causal discovery via NOTEARS with temporal forbidden masking ($H_2$), and fully neural graphical discovery-parameterization ($H_3$).

Across five paired random seeds under a 70/15/15 train/validation/test split, our empirical results establish three central findings:
(1) Granting scholarships ($T=1$) yields a statistically significant, positive causal average treatment effect on degree graduation ($ATE = +22.3\% \sim +23.7\%$, $p < 10^{-15}$), confirmed by DoWhy placebo and data-subset refutations.
(2) Neural CPD parameterization ($H_1, H_3$) markedly improves predictive probability calibration (lowering Brier score from $0.1719$ to $0.1667$) and discriminatory ranking (AUROC $= 0.8053$), while preserving structural causal invariance ($ATE \approx +0.23$).
(3) Causal mediation analysis demonstrates that $78.5\%$ of the scholarship effect operates via direct socioeconomic relief ($NDE = +16.96\%$), while $21.5\%$ is mediated through first-semester academic momentum ($NIE = +4.65\%$). 
Furthermore, Cinelli-Hazlett sensitivity analysis confirms robust unconfoundedness ($RV_{q=1} = 19.09\%$), and budget-constrained counterfactual policy learning reveals that prioritizing scholarship distribution based on Conditional Average Treatment Effects (CATE) produces $+164$ additional university graduates annually at zero incremental fiscal cost. Cross-dataset testing on 61,578 synthetic students validates that the learned graphical causal mechanisms remain invariant under distribution shifts ($ATE = +0.2143$, AUROC $= 0.8537$).
